\section{Discussion and Deployment Guideline}
\label{sec:discussion}

\subsection{What We Would Tell a Practitioner}
Four rules follow directly from the study. (1) \emph{Never ship a source-calibrated static threshold for a new domain.} Across four domain families the static threshold failed on six of six folds and in both directions; the failure is not a tuning problem but a distributional one. (2) \emph{Wrap the score with an online ACI layer and monitor its trace; use NA-ACI when anomalies may be abundant.} On five folds the layer recovered coverage within $0.015$ of nominal at negligible cost ($0.6$--$0.8\,\mu$s per decision); on anomaly-rich streams, NA-ACI keeps the budget on the normal regime and improves detection sharply at essentially unchanged FPR, at the price of one interpretable parameter $\kappa$ whose exclusion rate is monitorable online. (3) \emph{Report coverage and alarm rate separately, and know which one is guaranteed.} ACI controls the long-run alarm rate; normal-window FPR control follows when anomalies are rare relative to $\alphaNom$, and NA-ACI narrows the gap otherwise. (4) \emph{Audit the ranker before trusting the calibration.} Two folds here rank anomalies below chance; their coverage numbers look excellent and their alarms are worthless. A feature-baseline AUROC check (RMS or spectral energy) costs nothing and exposes this. A monitoring protocol closes the loop: (1) audit ranking with a feature-baseline AUROC, and stop if it is below chance; (2) run ACI by default, NA-ACI when prevalence is believed high; (3) watch three quantities---alarm-rate EWMA (band $[0.5\alphaNom,2\alphaNom]$), exclusion rate (compare against the prevalence estimate), and threshold drift (ADWIN/Page--Hinkley \cite{adwin2007,page1954})---with out-of-band signals triggering human inspection and persistent drift triggering recalibration; (4) recalibrate from a fresh labeled audit segment whenever the ranking audit fails.



\subsection{A Measurement Perspective}
Read through a metrology lens, the deployed threshold is a \emph{measured quantity}---the output of the calibration procedure---and this paper quantifies its transfer behavior: static calibration carries a transfer error above $0.7$ coverage on the worst fold, while the online layer holds it to an $O(\gamma)$ tracking error with an empirical seed spread below $0.02$. In VIM/GUM terms, the oracle decomposition of Section~\ref{sec:transfer} separates the \emph{domain term} (how far the target scale has moved) from the \emph{mechanism term} (how far the controller sits from the target-optimal threshold), and the coverage/alarm-rate pair is the accuracy/precision-style decomposition of the calibration, reported separately by construction. We stop short of a full uncertainty budget: threshold uncertainty is not propagated into event-level false-alarm probabilities, which is the natural next step for a measurement-audience treatment.

\subsection{Limitations}
The benchmark has six folds over four domain families, with single-dataset wind, one fab dataset, a photovoltaic family represented only by its synthetic-injection fold, and three of six folds targeting rotating machinery; "cross-domain" here means cross-population within industrial monitoring, not six separate industries. Pseudo-labels are unavoidable for several datasets, and PHM16's $q_{10}$ removal-rate protocol and the PVDAQ synthetic injection are the weakest links; both are documented and their sensitivity is reported (contamination sweep, diagnosis of below-chance folds). The encoder family is deliberately simple and uniform; we compare against mechanism-level reimplementations of the closest rivals (CICADA-, DACAD-style, Section~\ref{sec:rivals}) and against classical and foundation-model score producers (PCA, band-energy Mahalanobis, Chronos-bolt-small), while official-code comparisons remain a revision item. Finally, the conditional-coverage goal remains open: Mondrian and divergence-shifted variants did not deliver it, and NA-ACI narrows but does not close the conditional gap, so we claim no more than marginal coverage control with an improved budget allocation. Equation~\ref{eq:aci} is a scale-free variant of ACI; its long-run alarm-rate interpretation is $O(\gamma)$ in the log-threshold parameterization and does not inherit the exact finite-sample constants proved for the additive form. The source pools are also imbalanced ($1{,}024$--$16{,}007$ normal windows per dataset after capping, a $15.6\times$ spread); balanced-domain retraining remains untested. Dataset licences differ by provider (Paderborn via the KAt DataCenter terms, PVDAQ under NREL's terms, NASA and university datasets under research-use terms); we redistribute no raw data and defer licence specifics to each provider. The rival comparison is mechanism-level: our reimplementations preserve each rival's core mechanism (injection plus alignment head; expert mixture with divergence gating) under identical windows and training budgets, and official-code comparisons remain future work with the alignment protocol documented in the released repository.

\subsection{Future Work}
Three directions follow. First, conditional coverage beyond marginal control: NA-ACI fixes the allocation of the alarm budget but does not provide group-conditional guarantees, and the assumption behind its exclusion rule (anomaly-pure tail) deserves theory that links measurable diagnostics---exclusion rate, score-distribution drift---to coverage bounds. Second, extending the benchmark with more domain families and externally produced scores (foundation models, domain-adaptation detectors) to test whether the calibration conclusions and NA-ACI gains transfer to other score producers. Third, a theory-side account of threshold inflation: inflation spans two orders of magnitude across our folds and we can describe it but not yet predict it.
